This appendix holds the descriptive pass that precedes the Bayesian analysis. It covers the fifteen geometries of \autoref{tab:hfModels}, five to a figure, and reads each geometry twice: on the left what the model outputs, on the right what its internal state retains. Both readings average over phase, use a single signal family and carry no statement of uncertainty, so they are descriptive and give a high-level view of how frequency content behaves across the geometries. \autoref{fig:rec16-16} in \autoref{sec:results} is the $P=S=16$ row of the same measurement, drawn on its own.

\subsection{How the sweeps are produced}

The measurement is a sweep of pure sinusoids from $2$ to $250$\,Hz in one-hertz steps at $f_s=512$\,Hz, the same band used by the inferential stage, so that the band tasks of \autoref{tab:bandTasks} still partition what is measured. Each geometry receives a context of $512$ samples shortened to a whole number of strides, which gives $L=512$ for $S\in\{8,16,32\}$, $504$ for $S\in\{12,24\}$ and $500$ for $S=20$. An integer tone therefore completes a whole number of cycles in the context only for the first group, and the resulting leakage differs by geometry as a property of the measurement rather than of the model. As in \autoref{sec:geometries}, samples beyond the last full stride enter no token.

The left column is behavioural: the tone is given as context, the median forecast is appended to it and the model is run again, until $512$ samples have been generated; the dominant frequency of that window is the largest magnitude of its discrete Fourier transform after the mean is removed, averaged over four starting phases. The window length gives one-hertz bins for every geometry, so a dominant frequency is placed to within a hertz. The loop is imposed from outside, because Chronos-Bolt-Tiny emits its whole horizon in one step, so a curve that degrades along the sweep is partly the model reading its own output, with the error of each step carried into the next.

The right column is representational: at each swept frequency the internal state at the output quantile head is standardised, reduced to its leading twenty principal components, and a logistic regression is trained on that projection and scored by cross-validation, with the standardisation and the projection fitted inside each training fold. Eight tasks are scored per geometry: the seven hierarchical band tasks of \autoref{tab:bandTasks}, after Pagani et al.\autocite{paganiPreliminaryInsightsChronos2026}, which locate a tone by halving the analysed range three times over, and one local task separating a tone at $f-1$\,Hz from one at $f+1$\,Hz. A geometry that answers the first level and fails the third has kept the coarse position of a tone and lost its fine one. The band tasks are cross-validated with folds grouped by frequency, so a probe cannot score by memorising frequencies, while the local task holds out phases, using ten per frequency. Because the reduction is unsupervised and the probe is linear, what is scored is a lower bound on what the state holds: a red cell shows that a frequency is not easily readable, not that it is absent.

\begin{table}[!t]
    \centering
\small
    \setlength{\tabcolsep}{5pt}
    \begin{tabular}{@{}llrrr@{}}
        \toprule
        \textbf{Task} & \textbf{Level} & \textbf{Band [Hz]} & \textbf{Boundary [Hz]} & \textbf{Width [Hz]} \\
        \midrule
        \texttt{Mid} & 1 & $2$--$250$   & $126$ & $248$ \\
        \addlinespace[2pt]
        \texttt{L}   & 2 & $2$--$126$   & $64$  & $124$ \\
        \texttt{H}   & 2 & $127$--$250$ & $188$ & $123$ \\
        \addlinespace[2pt]
        \texttt{LL}  & 3 & $2$--$64$    & $33$  & $62$ \\
        \texttt{LH}  & 3 & $65$--$126$  & $95$  & $61$ \\
        \texttt{HL}  & 3 & $127$--$188$ & $157$ & $61$ \\
        \texttt{HH}  & 3 & $189$--$250$ & $219$ & $61$ \\
        \bottomrule
    \end{tabular}
\caption{The seven hierarchical band-classification tasks, after Pagani et al.\autocite{paganiPreliminaryInsightsChronos2026}, each with the level of split it belongs to, the band it is asked of, the boundary that divides that band and the width of the band. The boundaries partition $2$--$250$\,Hz and are held fixed across geometries, so a row may be read across the design.}
    \label{tab:bandTasks}
\end{table}


\begin{figure*}[!p]
    \centering
    \includegraphics[width=0.85\textwidth]{figures/FIG_E1_page1.png}
    \caption{Reconstruction and probe sweep, first of three: geometries p8-s8, p16-s8, p16-s12,
    p16-s16 and p24-s8. Each row is one geometry. Left, the dominant frequency of the generated
    window against the frequency supplied as context, averaged over four starting phases, drawn as
    the solid red curve. The pale red band around it is $\pm1$ standard deviation of that dominant
    frequency across the four phases, measured on the vertical axis and therefore in hertz: a wide
    band at a swept frequency means the frequency the model returns depends on where in its cycle
    the tone began, and a band that collapses to the curve means the reading is phase-independent
    there. The dotted diagonal marks a tone rebuilt correctly. Right, the
    accuracy of a linear probe on the internal state at each swept frequency, one strip per task,
    with the local $f\pm1$\,Hz task last. Dashed verticals are that geometry's own predicted
    candidates. This figure carries the two non-overlapping baselines, p8-s8 and p16-s16, beside the
    two geometries that share $S=8$ with different patch sizes.}
    \label{fig:sweepA}
\end{figure*}

\begin{figure*}[!p]
    \centering
    \includegraphics[width=0.85\textwidth]{figures/FIG_E1_page2.png}
    \caption{Second of three: geometries p24-s12, p24-s16, p24-s20, p24-s24 and p32-s8, read as \autoref{fig:sweepA}. The first four hold the patch size at $P=24$ and shorten the stride
    from $24$ to $12$, so the stride comb $c f_s/S$ moves across the band while the patch comb
    $k f_s/P$ stays where it is: this is the series on which the two branches of
    $\mathcal F_{\mathrm{lock}}$ can be told apart by eye, and p24-s16 and p24-s20 are among the
    six runs with $S \nmid P$ on which the branch estimates of \autoref{sec:models} rest.}
    \label{fig:sweepB}
\end{figure*}

\begin{figure*}[!p]
    \centering
    \includegraphics[width=0.85\textwidth]{figures/FIG_E1_page3.png}
    \caption{Third of three: geometries p32-s12, p32-s16, p32-s20, p32-s24 and p32-s32, read as \autoref{fig:sweepA}. The patch size is held at $P=32$, the widest of the design, so the
    patch comb is the densest and the overlap ratio runs from $0.62$ down to $0$ within this one figure.
    Read against \autoref{fig:sweepB} at equal stride, the difference between the two is a difference of
    patch size alone, which is the comparison the covariate $\log P$ of \autoref{eq:appM1:model}
    is fitted to.}
    \label{fig:sweepC}
\end{figure*}

\subsection{What is shown for each geometry}

Figures~\ref{fig:sweepA}, \ref{fig:sweepB} and~\ref{fig:sweepC} divide the design by patch size. The first carries the two smallest patches and the first $P=24$ run, where the combs are coarsest; the second completes the $P=24$ series and opens the $P=32$ block, the range over which stride varies at fixed patch size; the third holds the five $P=32$ geometries, where the patch comb is densest, ending on $P=S=32$, where the two branches coincide.

On the left, a point on the dotted diagonal is a tone rebuilt at the frequency it was given. A flat segment means that distinct inputs return dominant frequencies in the same bin, a loss of distinguishability rather than of amplitude, and a jump is a tone the model produces that was never supplied. The shaded band is one standard deviation over the starting phases. On the right, each cell is one swept frequency coloured by probe accuracy: green at or above $0.9$, amber in the middle, red below $0.6$, close to chance; grey marks a frequency no task covers. The dashed verticals on both sides are the candidate sites $\mathcal F_{\mathrm{lock}}$ of that geometry, computed from $(f_s,P,S)$ and not fitted. A frequency off the diagonal on the left but green on the right was lost between the representation and the output, not at the tokenisation, while a red cell means the internal state itself no longer separates that frequency from its neighbour.

\subsection{What the figures show}

Three things are visible on the sweep figures and hold across all fifteen geometries. The reconstruction curve follows the diagonal only over a low part of the band and departs from it well before the Nyquist limit; above that departure the returned dominant frequency is mostly a low-frequency value unrelated to the input, punctuated by isolated excursions. There is a broad plateau in the upper middle of the band on which many distinct inputs return one and the same dominant bin, which is the flat segment described above and is a loss of distinguishability rather than of amplitude. And the phase band is widest exactly on that plateau, so where the reading is worst it is also least stable across the starting phase.

The probe strips do not follow that pattern, since over most of the band, and for most of the eight tasks, the accuracy stays at or above $0.9$, including at frequencies the reconstruction no longer rebuilds correctly. Yellow and red cells occur, and cluster at the low end of the band and near some task boundaries, but they are the exception. Read across a row, the common case is therefore a frequency that the internal state still separates from its neighbour and that the generated output no longer places correctly, which locates the loss after the representation rather than at the tokenisation.

These readings are descriptive, since neither is paired against controls or carries an interval. What the figures show is that there is something to measure and that the two readings disagree in a specific direction: above the lower part of the band the output no longer places a tone that the internal state still separates. That observation gives the study its direction, pointing to where frequencies stop being represented in what the model returns, and it is the Bayesian hypotheses reported in \autoref{sec:bayesResults} that test it.
